% Options for packages loaded elsewhere
\PassOptionsToPackage{unicode}{hyperref}
\PassOptionsToPackage{hyphens}{url}
%
\documentclass[
]{article}
\usepackage{amsmath,amssymb}
\usepackage{iftex}
\ifPDFTeX
  \usepackage[T1]{fontenc}
  \usepackage[utf8]{inputenc}
  \usepackage{textcomp} % provide euro and other symbols
\else % if luatex or xetex
  \usepackage{unicode-math} % this also loads fontspec
  \defaultfontfeatures{Scale=MatchLowercase}
  \defaultfontfeatures[\rmfamily]{Ligatures=TeX,Scale=1}
\fi
\usepackage{lmodern}
\ifPDFTeX\else
  % xetex/luatex font selection
\fi
% Use upquote if available, for straight quotes in verbatim environments
\IfFileExists{upquote.sty}{\usepackage{upquote}}{}
\IfFileExists{microtype.sty}{% use microtype if available
  \usepackage[]{microtype}
  \UseMicrotypeSet[protrusion]{basicmath} % disable protrusion for tt fonts
}{}
\makeatletter
\@ifundefined{KOMAClassName}{% if non-KOMA class
  \IfFileExists{parskip.sty}{%
    \usepackage{parskip}
  }{% else
    \setlength{\parindent}{0pt}
    \setlength{\parskip}{6pt plus 2pt minus 1pt}}
}{% if KOMA class
  \KOMAoptions{parskip=half}}
\makeatother
\usepackage{xcolor}
\usepackage[margin=1in]{geometry}
\usepackage{longtable,booktabs,array}
\usepackage{calc} % for calculating minipage widths
% Correct order of tables after \paragraph or \subparagraph
\usepackage{etoolbox}
\makeatletter
\patchcmd\longtable{\par}{\if@noskipsec\mbox{}\fi\par}{}{}
\makeatother
% Allow footnotes in longtable head/foot
\IfFileExists{footnotehyper.sty}{\usepackage{footnotehyper}}{\usepackage{footnote}}
\makesavenoteenv{longtable}
\usepackage{graphicx}
\makeatletter
\def\maxwidth{\ifdim\Gin@nat@width>\linewidth\linewidth\else\Gin@nat@width\fi}
\def\maxheight{\ifdim\Gin@nat@height>\textheight\textheight\else\Gin@nat@height\fi}
\makeatother
% Scale images if necessary, so that they will not overflow the page
% margins by default, and it is still possible to overwrite the defaults
% using explicit options in \includegraphics[width, height, ...]{}
\setkeys{Gin}{width=\maxwidth,height=\maxheight,keepaspectratio}
% Set default figure placement to htbp
\makeatletter
\def\fps@figure{htbp}
\makeatother
\setlength{\emergencystretch}{3em} % prevent overfull lines
\providecommand{\tightlist}{%
  \setlength{\itemsep}{0pt}\setlength{\parskip}{0pt}}
\setcounter{secnumdepth}{-\maxdimen} % remove section numbering
\ifLuaTeX
  \usepackage{selnolig}  % disable illegal ligatures
\fi
\IfFileExists{bookmark.sty}{\usepackage{bookmark}}{\usepackage{hyperref}}
\IfFileExists{xurl.sty}{\usepackage{xurl}}{} % add URL line breaks if available
\urlstyle{same}
\hypersetup{
  hidelinks,
  pdfcreator={LaTeX via pandoc}}

\author{}
\date{\vspace{-2.5em}}

\begin{document}

\renewcommand{\arraystretch}{1.2}

Name: \vspace{0.1in}

\hrule

\hypertarget{m0409_quiz-v116}{%
\subsection{m0409\_quiz (v116)}\label{m0409_quiz-v116}}

\hypertarget{question-1}{%
\subsubsection{Question 1}\label{question-1}}

4 candidates (W through Z) are running in an election with a single
winner. There are 8 voters (A through H) who each rank the candidates
from 1 to 4, with 1 meaning first choice and 4 meaning last choice.
Those rankings are shown below. Each column represents a single voter's
ranking.

\begin{longtable}[]{@{}
  >{\centering\arraybackslash}p{(\columnwidth - 16\tabcolsep) * \real{0.1111}}
  >{\centering\arraybackslash}p{(\columnwidth - 16\tabcolsep) * \real{0.1111}}
  >{\centering\arraybackslash}p{(\columnwidth - 16\tabcolsep) * \real{0.1111}}
  >{\centering\arraybackslash}p{(\columnwidth - 16\tabcolsep) * \real{0.1111}}
  >{\centering\arraybackslash}p{(\columnwidth - 16\tabcolsep) * \real{0.1111}}
  >{\centering\arraybackslash}p{(\columnwidth - 16\tabcolsep) * \real{0.1111}}
  >{\centering\arraybackslash}p{(\columnwidth - 16\tabcolsep) * \real{0.1111}}
  >{\centering\arraybackslash}p{(\columnwidth - 16\tabcolsep) * \real{0.1111}}
  >{\centering\arraybackslash}p{(\columnwidth - 16\tabcolsep) * \real{0.1111}}@{}}
\toprule\noalign{}
\begin{minipage}[b]{\linewidth}\centering
\end{minipage} & \begin{minipage}[b]{\linewidth}\centering
voter A
\end{minipage} & \begin{minipage}[b]{\linewidth}\centering
voter B
\end{minipage} & \begin{minipage}[b]{\linewidth}\centering
voter C
\end{minipage} & \begin{minipage}[b]{\linewidth}\centering
voter D
\end{minipage} & \begin{minipage}[b]{\linewidth}\centering
voter E
\end{minipage} & \begin{minipage}[b]{\linewidth}\centering
voter F
\end{minipage} & \begin{minipage}[b]{\linewidth}\centering
voter G
\end{minipage} & \begin{minipage}[b]{\linewidth}\centering
voter H
\end{minipage} \\
\midrule\noalign{}
\endhead
\bottomrule\noalign{}
\endlastfoot
candidate W & 3 & 2 & 1 & 4 & 4 & 1 & 4 & 1 \\
candidate X & 4 & 3 & 2 & 1 & 1 & 2 & 2 & 4 \\
candidate Y & 2 & 4 & 3 & 2 & 3 & 4 & 3 & 3 \\
candidate Z & 1 & 1 & 4 & 3 & 2 & 3 & 1 & 2 \\
\end{longtable}

\begin{enumerate}
\def\labelenumi{\arabic{enumi}.}
\tightlist
\item
  What \href{https://en.wikipedia.org/wiki/Borda_count}{Borda count}
  would each candidate get?
\item
  Who wins a Borda-count election?
\end{enumerate}

\color{red}

12 13 8 15 \ldots{} Z

\color{black}

\vfill

\hypertarget{question-2}{%
\subsubsection{Question 2}\label{question-2}}

5 candidates (V through Z) are running in an election with a single
winner. There are 11 voters (A through K) who each rank the candidates
from 1 to 5, with 1 meaning first choice and 5 meaning last choice.
Those rankings are shown below. Each column represents a single voter's
ranking.


\includegraphics{m0409/m0409_SOL_116_files/figure-latex/unnamed-chunk-6-1.pdf}\\

Who would win an
\href{https://en.wikipedia.org/wiki/Instant-runoff_voting}{instant-runoff
election}?

\color{red}

Candidate Y wins!

\color{black}

\hfill\break

\newpage

\hypertarget{question-3}{%
\subsubsection{Question 3}\label{question-3}}

4 candidates (W through Z) are running in an election with a single
winner. There are 7 voters (A through G) who each rank the candidates
from 1 to 4, with 1 meaning first choice and 4 meaning last choice.
Those rankings are shown below. Each column represents a single voter's
ranking.

\includegraphics{m0409/m0409_SOL_116_files/figure-latex/unnamed-chunk-9-1.pdf}\\

A \href{https://en.wikipedia.org/wiki/Condorcet_method}{Condorcet
winner} would win in a head-to-head election against any other
candidate. Who is the Condorcet winner? (Or, possibly, there is none.)

\color{red}

Notice Y wins against all other candidates. Thus Y is the Condorcet
winner.

\color{black}

\hypertarget{question-4}{%
\subsubsection{Question 4}\label{question-4}}

In an \href{https://en.wikipedia.org/wiki/Approval_voting}{approval-vote
election}, there are 4 candidates and 73 voters. Each voter indicated
which candidates they approve of (with a ``Y'') or disapprove of (with a
``N''). All possible ballots are shown below, and the table indicates
how many of each type of ballot was cast.


\includegraphics{m0409/m0409_SOL_116_files/figure-latex/unnamed-chunk-12-1.pdf}\\

\begin{enumerate}
\def\labelenumi{\arabic{enumi}.}
\tightlist
\item
  How many approvals did Candidate W get?
\item
  How many approvals did Candidate X get?
\item
  How many approvals did Candidate Y get?
\item
  How many approvals did Candidate Z get?
\item
  Who won the election?
\end{enumerate}

\color{red}

29 45 32 50 \ldots{} So Z wins.

\color{black}

\end{document}
